\section{Full result tables}
\label{app:tables}

This appendix collects the tables that the preceding appendices read against.
Every table is generated by the workflow from the same caches as the figures in the main text.

\tabref{encoder-matrix} repeats the field tasks for the three base models of Appendix~\ref{app:benchmarks}.
\tabref{radial-vec2text} is the \texttt{vec2text} magnitude sweep of Appendix~\ref{app:length-dial}.
\tabref{hierarchy-labels} closes the appendix with the native output of every embedding method for every node of the PACS slice scored in Appendix~\ref{app:baseline-labels}.

\begin{table}[htbp]
\centering
\small
\caption{Per-encoder field-task matrix. Each cell reports the two \doctolora variants as raw\,/\,+$g_\theta$ (the frozen embedding and the single general citation transform), with metrics, evaluation pools, and bootstrap protocol as in \tabref{similarity} (standard deviations omitted). Text-baseline scores are encoder-independent and given in \tabref{similarity}.}
\label{tab:encoder-matrix}
% Auto-generated by workflow/plot/tab_similarity_benchmarks.py -- do not edit by hand.
\begin{tabular}{llccc}
\toprule
Encoder & Field & Next-paper (AUC) & Topic (F1) & Collaboration (AUC) \\
\midrule
\multirow{3}{*}{Gemma-2-2b}
 & Economics & .792\,/\,.888 & .217\,/\,.317 & .592\,/\,.617 \\
 & Psychology & .786\,/\,.910 & .180\,/\,.319 & .593\,/\,.673 \\
 & Physics & .864\,/\,.946 & .569\,/\,.582 & .827\,/\,.823 \\
\midrule
\multirow{3}{*}{Qwen3-4B}
 & Economics & .813\,/\,.910 & .242\,/\,.374 & .596\,/\,.629 \\
 & Psychology & .799\,/\,.923 & .217\,/\,.340 & .602\,/\,.687 \\
 & Physics & .880\,/\,.958 & .590\,/\,.590 & .792\,/\,.795 \\
\midrule
\multirow{3}{*}{Mistral-7B}
 & Economics & .825\,/\,.913 & .234\,/\,.368 & .599\,/\,.639 \\
 & Psychology & .815\,/\,.926 & .217\,/\,.348 & .601\,/\,.681 \\
 & Physics & .876\,/\,.965 & .591\,/\,.595 & .820\,/\,.781 \\
\bottomrule
\end{tabular}

\end{table}



\begin{table}[htbp]
\centering
\footnotesize
\setlength{\tabcolsep}{3pt}
\caption{\texttt{vec2text} magnitude sweep on the same five seed documents. \texttt{vec2text} is an inverter and takes no prompt, so we read back its reconstruction directly. The reconstruction stays on topic at the same surface specificity for larger $\alpha$ and degrades into off-topic fragments at small $\alpha$, never abstracting toward the broader field. Entries truncated to leading terms.}
\label{tab:radial-vec2text}
\rowcolors{2}{gray!12}{white}
\begin{tabularx}{\linewidth}{@{}c *{5}{>{\raggedright\arraybackslash}X}@{}}
\toprule
$\alpha$ & \textbf{Physics} & \textbf{Biochemistry} & \textbf{Fruit} & \textbf{Person} & \textbf{Patent (LED)} \\
\midrule
$1.15$ & is the fractional quantum effect of accumulating electrons in certain atoms.\,\dots & (abbreviated CY), commonly ATP, is an enzyme cycle for the production\,\dots & a banana is an edible plantation fruit, traditionally specializing in long,\,\dots & Nikola Tesla, an American engineer from Serbia, was an 1860s American\,\dots & comprised an ion-doped light-emitting semiconductor forming a nitride structure, nitrid \\
$0.9$ & quantum fraction is a function of the Hall effect in physics:\,\dots & enzyme cycle is the metabolic use of ATP for generating cryptic\,\dots & fruit is a elongated banana. This is due to the variety\,\dots & Nikola Tesla was a Serbian engineer and American engineer, known for\,\dots & structure comprised a light-emitting compound of a double-type nitride semiconductor (e.g. \\
$0.7$ & quantum mechanics is the fraction of the Hall effect which is\,\dots & acids used for the production of glycogen. The Krebs cycle is\,\dots & plant species is a banana. The name derives from a linguistic\,\dots & engineering genius Nikola Tesla is a Serbian engineer who contributed to\,\dots & structure of a double-layered nickel semiconductor (indicated in the introduction of\,\dots \\
$0.55$ & equation for the physics of fractional heating is a contribution to\,\dots & the CYCY process of carbohydrates is a way of generating ATP,\,\dots & botanical designation of a plant. A banana is a combination of\,\dots & American mechanical engineer. Nikola Tesla is credited for the production of\,\dots & layer of a g-type nickel plating (indicated in the following table:\,\dots \\
$0.42$ & equation for quantum mechanics is a result of the publication of\,\dots & the metabolic cycle of enzymes is a contribution to the aforementioned\,\dots & botanical name for a banana. This is a variant of the\,\dots & American electrical engineer. Nikola Tesla is credited for the contribution of\,\dots & a GLTI semiconductor pairing of the following structures: (a) (initiated) (a)\,\dots \\
$0.3$ & a number of awards in the field of physics & biology is a key component of the Crypt of the Chemical\,\dots & biology of plants is a classic in the discipline of banana\,\dots & American electrical engineer who was a major contributor to the creation\,\dots & a layered semiconductor for the neo-Neo-Geo-Neo-Neo-Dai \\
$0.2$ & a prestigious award in the field of physics for the 1860\,\dots & biology is a classic in the vein of the MIT CYPATHE,\,\dots & classic in Asian agriculture. The title of a chapter is a\,\dots & American mechanical engineer. Tesla contributed to a number of awards including\,\dots & a DP-type semiconductor engineering is \\
$0.12$ & a prestigious journal in the area of physics entitled & a.k.a. The Chemical Cycle of the As - a -does it\,\dots & classic in the history of the Venus Project & a history of Tesla engineering. Contributor to the 97th edition of\,\dots & a fusion of the following disciplines \\
$0.06$ & classic in the history of physics. 109. & a.k.a. 179. 2004: The Academy of Sciences published a review of\,\dots & classic in the history of engineering. William Rice Publishers & 1925. The publication of a book on the history of electricity\,\dots & a.k.a. DP Engineering \\
\bottomrule
\end{tabularx}
\end{table}

\begin{footnotesize}
\setlength{\tabcolsep}{3pt}
\renewcommand{\arraystretch}{1.2}
\rowcolors{3}{gray!12}{white}
\begin{longtable}{@{}>{\raggedright\arraybackslash}p{1.85cm} >{\raggedright\arraybackslash}p{2.55cm} >{\raggedright\arraybackslash}p{2.6cm} >{\raggedright\arraybackslash}p{2.55cm} >{\raggedright\arraybackslash}p{2.55cm}@{}}
\caption{
Cluster labels for every node of a four-branch PACS slice. Each cell is the method's native output, exactly the string scored in \tabref{label-eval}, truncated at a word boundary for the page but scored in full. \doctolora and \texttt{ICAE} answer the field prompt; \texttt{KeyLLM} returns a keyword list and \texttt{vec2text} a reconstruction, which is why only the first two columns read as field names. Generated from the evaluation's own table by \texttt{make\_hierarchy\_rows.py}, so it cannot drift from \tabref{label-eval}. The first column gives the PACS code and a short name, indented by tree depth, with field rows in bold.}
\label{tab:hierarchy-labels} \\
\toprule
\rowcolor{white}
& \multicolumn{4}{c}{\textbf{Native output (as scored)}} \\
\cmidrule(lr){2-5}
\rowcolor{white}
\textbf{PACS node} & \textbf{\doctolora} & \textbf{\texttt{KeyLLM}} & \textbf{\texttt{ICAE}} & \textbf{\texttt{vec2text}} \\
\midrule
\endfirsthead
\toprule
\rowcolor{white}
& \multicolumn{4}{c}{\textbf{Native output (as scored)}} \\
\cmidrule(lr){2-5}
\rowcolor{white}
\textbf{PACS node} & \textbf{\doctolora} & \textbf{\texttt{KeyLLM}} & \textbf{\texttt{ICAE}} & \textbf{\texttt{vec2text}} \\
\midrule
\endhead
\bottomrule
\endlastfoot
\textbf{0~~General} & \textbf{Quantum field theory} & partial decoherence, thermalization, time-domain\,\dots & Quantum many-body systems & coupled-dissonance correlations in one-way systems. The\,\dots \\
\hspace{0.8em}03~~Quantum mechanics, QFT & \textbf{Quantum optics} & quantum state transfer, entanglement, four qubits\,\dots & Quantum optics & coupled entanglement of one-state quantums, where the\,\dots \\
\hspace{1.8em}03.65~~Quantum mechanics & \textbf{Quantum mechanics} & quantum decoherence, single qubit, random matrix theory\,\dots & Quantum mechanics & as an empirical measure of the interpolation of one-state\,\dots \\
\hspace{1.8em}03.67~~Quantum information & \textbf{Quantum information} & Measurement-based quantum computing, AKLT state, spin-1\,\dots & Quantum information & entanglement of single-state quantum computations, so\,\dots \\
\hspace{0.8em}05~~Statistical physics & \textbf{Statistical physics} & Disturbance spreading, Incommensurate systems\,\dots & Statistical mechanics & diffuse-phase dynamics in univariate quantization. This\,\dots \\
\hspace{1.8em}05.40~~Fluctuations, noise & \textbf{Statistical physics} & bistable systems, colored noise, stochastic relaxation\,\dots & Statistical physics & diffusion-display coefficients in one-variable\,\dots \\
\hspace{1.8em}05.45~~Nonlinear dynamics, chaos & \textbf{Nonlinear dynamics} & separatrix map, resonant dynamics, global chaos onset\,\dots & Nonlinear dynamics & coupled chaotic quantization in classical dynamics. This\,\dots \\
\addlinespace[2pt]
\textbf{2~~Classical physics} & \textbf{Nonlinear optics} & Nonlinear light pulse propagation, Polarized light\,\dots & Quantum optics & multi-mode scattering of neutrons via classical optical\,\dots \\
\hspace{0.8em}42~~Optics & \textbf{Quantum optics} & Photon antibunching, Photonic molecule, Coupled cavities\,\dots & Quantum optics & coupled photon-diametric scattering of neutrons in\,\dots \\
\hspace{1.8em}42.50~~Quantum optics & \textbf{Quantum optics} & mesoscopic states, cavity-field state, atom-field\,\dots & Quantum optics & two-level physics involving scattering of coherent atoms\,\dots \\
\hspace{1.8em}42.65~~Nonlinear optics & \textbf{Nonlinear optics} & solitons, spontaneous symmetry breaking, photonic\,\dots & Nonlinear optics & nonsimultaneous low-wave propagation in coherent optical\,\dots \\
\hspace{0.8em}47~~Fluid dynamics & \textbf{Fluid dynamics} & Two-dimensional turbulence, freely decaying turbulence\,\dots & Fluid dynamics & fluid-displacement dynamics in experimentally correlated\,\dots \\
\hspace{1.8em}47.20~~Flow instabilities & \textbf{Fluid dynamics} & Bifurcation diversity, Thermocapillary liquid layers\,\dots & Fluid dynamics & in-convection instability in linear-wavelength fluid\,\dots \\
\hspace{1.8em}47.27~~Turbulent flows & \textbf{Fluid dynamics} & Mean-field approximation, turbulence theory\,\dots & Turbulence theory & in turbulence-displacement dynamics, the initial\,\dots \\
\addlinespace[2pt]
\textbf{3~~Condensed matter} & \textbf{Condensed matter physics} & Quasi-one-dimensional electron systems, t-J model\,\dots & Condensed matter physics & in-distance interactions between electron-bonding\,\dots \\
\hspace{0.8em}71~~Electronic structure & \textbf{Condensed matter physics} & semiclassical action, dynamical mean-field theory, local\,\dots & Condensed matter physics & phase-dimersion configuration of classical electron-fluid\,\dots \\
\hspace{1.8em}71.10~~Many-electron systems & \textbf{Condensed matter physics} & Mott physics, Mott transition, spin fluctuations\,\dots & Condensed Matter Physics & one-dimensional quasi-dotonal model where the Hubbard\,\dots \\
\hspace{1.8em}71.20~~Bulk band structure & \textbf{Condensed matter physics} & Electronic structure, Pyrochlore metals, Cd2Os2O7\,\dots & Materials science & (dBa)-structure calculations of the respective metallic\,\dots \\
\hspace{0.8em}75~~Magnetic materials & \textbf{Condensed matter physics} & Quantum Monte Carlo, sign problem, itinerant\,\dots & Condensed Matter Physics & inferromagnetic spin-diabetic scattering interactions\,\dots \\
\hspace{1.8em}75.10~~Magnetic ordering & \textbf{Condensed matter physics} & square-lattice Heisenberg antiferromagnet, S=1/2\,\dots & Condensed Matter Physics & in anastomosing-phase-spatial-inverse quantum state, with\,\dots \\
\hspace{1.8em}75.30~~Ordered magnets & \textbf{Condensed matter physics} & metal-insulator transition, La1-xCaxMnO3, manganites\,\dots & Magnetism & inferomorphic magnetic spin-phases at 0–0,0–0,0–0,0–0,0–0 \\
\addlinespace[2pt]
\textbf{6~~Elementary particles} & \textbf{Particle physics} & Left-right symmetry, LHC, SU(2)\_L x SU(2)\_R x U(1)\_B-L\,\dots & High Energy Physics & scattering of resonant singularity quarks on the\,\dots \\
\hspace{0.8em}11~~Fields and particles & \textbf{Quantum field theory} & light-front Hamiltonian, spontaneous chiral symmetry\,\dots & Quantum Field Theory & with the symmetries of fermion quantum diffusion. This\,\dots \\
\hspace{1.8em}11.10~~Field theory & \textbf{Quantum field theory} & sine-Gordon field theory, massive, (1+1)-dimensional\,\dots & Quantum field theory & symmetries of the classical quantization-dynamic theory\,\dots \\
\hspace{1.8em}11.15~~Gauge field theories & \textbf{Quantum field theory} & SU(2) gauge theory, adjoint Dirac flavor, infrared\,\dots & Quantum field theory & scattering of d-quantity quantums. This d-quantity\,\dots \\
\hspace{0.8em}12~~Particle systematics & \textbf{Particle physics} & Left-Right symmetry, LHC, SU(2)\_L x SU(2)\_R x U(1)\_B-L\,\dots & High Energy Physics & scattering of low-mass quark quantities. This approach\,\dots \\
\hspace{1.8em}12.38~~Quantum chromodynamics & \textbf{Quantum chromodynamics} & Nucleon strange quark content, Lattice QCD, Two-flavor\,\dots & Nuclear physics & with d-quantity scattering of d-quantity laton particles\,\dots \\
\hspace{1.8em}12.60~~Beyond standard model & \textbf{Particle physics} & SO(10), supersymmetry breaking, gauge mediation, grand\,\dots & Particle physics & high scattering-symmetry sbN, the experimental model will\,\dots \\

\end{longtable}
\end{footnotesize}
